% Options for packages loaded elsewhere
\PassOptionsToPackage{unicode}{hyperref}
\PassOptionsToPackage{hyphens}{url}
\documentclass[
  12pt,
]{ctexart}
\usepackage{xcolor}
\usepackage{amsmath,amssymb}
\setcounter{secnumdepth}{-\maxdimen} % remove section numbering
\usepackage{iftex}
\ifPDFTeX
  \usepackage[T1]{fontenc}
  \usepackage[utf8]{inputenc}
  \usepackage{textcomp} % provide euro and other symbols
\else % if luatex or xetex
  \usepackage{unicode-math} % this also loads fontspec
  \defaultfontfeatures{Scale=MatchLowercase}
  \defaultfontfeatures[\rmfamily]{Ligatures=TeX,Scale=1}
\fi
\usepackage{lmodern}
\ifPDFTeX\else
  % xetex/luatex font selection
\fi
% Use upquote if available, for straight quotes in verbatim environments
\IfFileExists{upquote.sty}{\usepackage{upquote}}{}
\IfFileExists{microtype.sty}{% use microtype if available
  \usepackage[]{microtype}
  \UseMicrotypeSet[protrusion]{basicmath} % disable protrusion for tt fonts
}{}
\makeatletter
\@ifundefined{KOMAClassName}{% if non-KOMA class
  \IfFileExists{parskip.sty}{%
    \usepackage{parskip}
  }{% else
    \setlength{\parindent}{0pt}
    \setlength{\parskip}{6pt plus 2pt minus 1pt}}
}{% if KOMA class
  \KOMAoptions{parskip=half}}
\makeatother
\usepackage{color}
\usepackage{fancyvrb}
\newcommand{\VerbBar}{|}
\newcommand{\VERB}{\Verb[commandchars=\\\{\}]}
\DefineVerbatimEnvironment{Highlighting}{Verbatim}{commandchars=\\\{\}}
% Add ',fontsize=\small' for more characters per line
\newenvironment{Shaded}{}{}
\newcommand{\AlertTok}[1]{\textcolor[rgb]{1.00,0.00,0.00}{\textbf{#1}}}
\newcommand{\AnnotationTok}[1]{\textcolor[rgb]{0.38,0.63,0.69}{\textbf{\textit{#1}}}}
\newcommand{\AttributeTok}[1]{\textcolor[rgb]{0.49,0.56,0.16}{#1}}
\newcommand{\BaseNTok}[1]{\textcolor[rgb]{0.25,0.63,0.44}{#1}}
\newcommand{\BuiltInTok}[1]{\textcolor[rgb]{0.00,0.50,0.00}{#1}}
\newcommand{\CharTok}[1]{\textcolor[rgb]{0.25,0.44,0.63}{#1}}
\newcommand{\CommentTok}[1]{\textcolor[rgb]{0.38,0.63,0.69}{\textit{#1}}}
\newcommand{\CommentVarTok}[1]{\textcolor[rgb]{0.38,0.63,0.69}{\textbf{\textit{#1}}}}
\newcommand{\ConstantTok}[1]{\textcolor[rgb]{0.53,0.00,0.00}{#1}}
\newcommand{\ControlFlowTok}[1]{\textcolor[rgb]{0.00,0.44,0.13}{\textbf{#1}}}
\newcommand{\DataTypeTok}[1]{\textcolor[rgb]{0.56,0.13,0.00}{#1}}
\newcommand{\DecValTok}[1]{\textcolor[rgb]{0.25,0.63,0.44}{#1}}
\newcommand{\DocumentationTok}[1]{\textcolor[rgb]{0.73,0.13,0.13}{\textit{#1}}}
\newcommand{\ErrorTok}[1]{\textcolor[rgb]{1.00,0.00,0.00}{\textbf{#1}}}
\newcommand{\ExtensionTok}[1]{#1}
\newcommand{\FloatTok}[1]{\textcolor[rgb]{0.25,0.63,0.44}{#1}}
\newcommand{\FunctionTok}[1]{\textcolor[rgb]{0.02,0.16,0.49}{#1}}
\newcommand{\ImportTok}[1]{\textcolor[rgb]{0.00,0.50,0.00}{\textbf{#1}}}
\newcommand{\InformationTok}[1]{\textcolor[rgb]{0.38,0.63,0.69}{\textbf{\textit{#1}}}}
\newcommand{\KeywordTok}[1]{\textcolor[rgb]{0.00,0.44,0.13}{\textbf{#1}}}
\newcommand{\NormalTok}[1]{#1}
\newcommand{\OperatorTok}[1]{\textcolor[rgb]{0.40,0.40,0.40}{#1}}
\newcommand{\OtherTok}[1]{\textcolor[rgb]{0.00,0.44,0.13}{#1}}
\newcommand{\PreprocessorTok}[1]{\textcolor[rgb]{0.74,0.48,0.00}{#1}}
\newcommand{\RegionMarkerTok}[1]{#1}
\newcommand{\SpecialCharTok}[1]{\textcolor[rgb]{0.25,0.44,0.63}{#1}}
\newcommand{\SpecialStringTok}[1]{\textcolor[rgb]{0.73,0.40,0.53}{#1}}
\newcommand{\StringTok}[1]{\textcolor[rgb]{0.25,0.44,0.63}{#1}}
\newcommand{\VariableTok}[1]{\textcolor[rgb]{0.10,0.09,0.49}{#1}}
\newcommand{\VerbatimStringTok}[1]{\textcolor[rgb]{0.25,0.44,0.63}{#1}}
\newcommand{\WarningTok}[1]{\textcolor[rgb]{0.38,0.63,0.69}{\textbf{\textit{#1}}}}
\usepackage{longtable,booktabs,array}
\usepackage{calc} % for calculating minipage widths
% Correct order of tables after \paragraph or \subparagraph
\usepackage{etoolbox}
\makeatletter
\patchcmd\longtable{\par}{\if@noskipsec\mbox{}\fi\par}{}{}
\makeatother
% Allow footnotes in longtable head/foot
\IfFileExists{footnotehyper.sty}{\usepackage{footnotehyper}}{\usepackage{footnote}}
\makesavenoteenv{longtable}
\setlength{\emergencystretch}{3em} % prevent overfull lines
\providecommand{\tightlist}{%
  \setlength{\itemsep}{0pt}\setlength{\parskip}{0pt}}
\usepackage{bookmark}
\IfFileExists{xurl.sty}{\usepackage{xurl}}{} % add URL line breaks if available
\urlstyle{same}
\hypersetup{
  hidelinks,
  pdfcreator={LaTeX via pandoc}}

\author{SCX}
\date{2026}

\begin{document}

\section{SCX
两层描述符：错误驱动的状态发现}\label{scx-ux4e24ux5c42ux63cfux8ff0ux7b26ux9519ux8befux9a71ux52a8ux7684ux72b6ux6001ux53d1ux73b0}

\begin{quote}
SCX 的核心理念：状态不由人来定义，而由''模型在哪里失败''来定义。

这是 SCX 区别于所有基于人类特征聚类方法的本质差异。
\end{quote}

\textbf{文档定位}：Paper 2 (Residual-State Maps) 与 Paper 4 (SCX Theory)
之间的桥梁。 Paper 2
证明了''误差有结构根源''，两层描述符就是发现这些根源的方法； Paper 4
的''状态条件专家性''依赖两层描述符来定义什么是''状态''。

\textbf{依赖的前序文档}： - \texttt{01\_SCX\_核心框架\_数学分析.md} ---
状态条件风险的形式化 - \texttt{05\_数学根源与证明.md} ---
信息瓶颈与状态划分的测度论基础 - \texttt{06\_实现架构与实验设计.md} ---
SCX 实现架构

\textbf{相关实现}： - \texttt{src/scx/encoders/error\_driven.py} ---
ErrorDrivenEncoder - \texttt{src/scx/state/two\_layer.py} ---
TwoLayerStateDiscovery - \texttt{src/scx/encoders/base.py} ---
SCXStateEncoder 基类 - \texttt{src/scx/state/discovery.py} ---
StateDiscovery 基类

\begin{center}\rule{0.5\linewidth}{0.5pt}\end{center}

\subsection{目录}\label{ux76eeux5f55}

\begin{enumerate}
\def\labelenumi{\arabic{enumi}.}
\tightlist
\item
  \hyperref[1-ux95eeux9898ux4e3aux4ec0ux4e48ux4e00ux5c42ux63cfux8ff0ux7b26ux4e0dux591f]{问题：为什么一层描述符不够？}
\item
  \hyperref[2-ux4e24ux5c42ux63cfux8ff0ux7b26ux4f53ux7cfb]{两层描述符体系}
\item
  \hyperref[3-ux6570ux5b66ux5f62ux5f0fux5316]{数学形式化}
\item
  \hyperref[4-scx-ux5b9eux73b0]{SCX 实现}
\item
  \hyperref[5-ux4e0eux7adeux4e89ux65b9ux6cd5ux7684ux5bf9ux6bd4]{与竞争方法的对比}
\item
  \hyperref[6-ux5b9eux9a8cux9a8cux8bc1]{实验验证}
\item
  \hyperref[7-ux8bbaux6587ux5b9aux4f4dux4e0eux8d21ux732e]{论文定位与贡献}
\item
  \hyperref[8-ux53c2ux8003ux6587ux732e]{参考文献}
\end{enumerate}

\begin{center}\rule{0.5\linewidth}{0.5pt}\end{center}

\subsection{1.
问题：为什么一层描述符不够？}\label{ux95eeux9898ux4e3aux4ec0ux4e48ux4e00ux5c42ux63cfux8ff0ux7b26ux4e0dux591f}

\subsubsection{1.1
人类描述符的局限性}\label{ux4ebaux7c7bux63cfux8ff0ux7b26ux7684ux5c40ux9650ux6027}

现有势函数开发的共同路径：研究人员凭直觉选择描述符（元素种类、晶体结构、温度范围），然后在这些描述符定义的类别上做均匀采样。这一范式的根本缺陷在于：

\textbf{人类选的描述符可能不捕捉模型的真实失败模式。}

\paragraph{1.1.1 AlN v3 案例}\label{aln-v3-ux6848ux4f8b}

在 AlN v3 的实际开发中，这一问题暴露得极为清晰：

\begin{itemize}
\tightlist
\item
  研究人员定义了 12
  维手工特征（化学组分、结构类型、温度区段等），试图覆盖所有的化学环境
\item
  但使用这些特征对测试集做聚类后，约 \textbf{50\% 的帧被挤在一个状态中}
\item
  这个''超级状态''内部包含着物理上完全不同的结构（表面、缺陷、高应变构型），它们的力误差差异可达一个数量级
\item
  SCX 在这 50\%
  帧的区域内完全失效------因为单一状态内误差方差太大，任何全局专家路由策略都无法有效区分
\end{itemize}

根源分析：人类描述符的维度选择受限于研究人员的\textbf{认知盲区}------人只能注意到自己已知的结构类别，无法预见模型会在哪些''意想不到''的构型中失效。

\paragraph{1.1.2 Layer 1
编码的信息瓶颈}\label{layer-1-ux7f16ux7801ux7684ux4fe1ux606fux74f6ux9888}

形式化地，设 Layer 1 编码为
\(\varphi_1: \mathcal{X} \rightarrow \mathbb{R}^d\)，残差为
\(r(x) = |f_{\text{model}}(x) - f_{\text{true}}(x)|\)。一层编码的信息瓶颈是：

\[I(\varphi_1(X); r) \leq I(X; r)\]

如果 \(\varphi_1\) 的维度选择与 \(r\)
不匹配，不等式左边远小于右边。换言之，\(\varphi_1\)
丢弃了预测残差中的关键信息。

\textbf{更根本的问题}：\(\varphi_1\) 的编码空间 \(\mathbb{R}^d\)
由人类先验定义，它优先编码''人认为重要的结构特征''（如晶格常数、对称性），而非''模型失败的模式''。优化
\(I(\varphi_1; r)\) 与优化 \(I(\varphi_1; X)\)
是不同目标------而人类描述符天然偏向后者。

\paragraph{1.1.3 为什么 PCA/SVD
不能解决}\label{ux4e3aux4ec0ux4e48-pcasvd-ux4e0dux80fdux89e3ux51b3}

一个自然的反驳：既然 12 维嵌入导致 50\% 帧拥挤，为什么不对 \(\varphi_1\)
做 PCA 降维或增加维度？

\begin{itemize}
\tightlist
\item
  \textbf{增加维度}：\(d=12\)
  已是手工特征的极限。继续增加维度面临维度灾难------更多维意味着更稀疏的采样，且人类无法可靠地定义更多有物理意义的描述符
\item
  \textbf{PCA 降维}：PCA 优化的是
  \(\text{Var}(\varphi_1)\)（方差最大方向），而非
  \(I(\varphi_1; r)\)（与误差的相关性）。PCA
  的前几个主成分几乎总是描述全局结构差异（如晶胞体积变化），而非局域的误差模式
\item
  \textbf{UMAP/t-SNE}：可视化友好，但保留的是局部邻域结构，仍不直接针对
  \(r\) 优化
\end{itemize}

\textbf{关键洞察}：问题不在于''维度不够''或''空间不平衡''，而在于\textbf{目标错配}------现有描述符的设计目标是''完整描述结构''，而非''解释误差''。

\subsubsection{1.2
竞争方法的共同缺陷}\label{ux7adeux4e89ux65b9ux6cd5ux7684ux5171ux540cux7f3aux9677}

\begin{longtable}[]{@{}
  >{\raggedright\arraybackslash}p{(\linewidth - 8\tabcolsep) * \real{0.1200}}
  >{\raggedright\arraybackslash}p{(\linewidth - 8\tabcolsep) * \real{0.2200}}
  >{\raggedright\arraybackslash}p{(\linewidth - 8\tabcolsep) * \real{0.2200}}
  >{\raggedright\arraybackslash}p{(\linewidth - 8\tabcolsep) * \real{0.2600}}
  >{\raggedright\arraybackslash}p{(\linewidth - 8\tabcolsep) * \real{0.1800}}@{}}
\toprule\noalign{}
\begin{minipage}[b]{\linewidth}\raggedright
方法
\end{minipage} & \begin{minipage}[b]{\linewidth}\raggedright
状态定义者
\end{minipage} & \begin{minipage}[b]{\linewidth}\raggedright
描述符来源
\end{minipage} & \begin{minipage}[b]{\linewidth}\raggedright
是否误差驱动
\end{minipage} & \begin{minipage}[b]{\linewidth}\raggedright
核心缺陷
\end{minipage} \\
\midrule\noalign{}
\endhead
\bottomrule\noalign{}
\endlastfoot
DIRECT sampling & 人 (SOAP features) & 人类预设 & ❌ & SOAP
描述结构而非误差；聚类边界与模型失败模式无关 \\
DP-GEN & 人 (model deviation) & model deviation 标量 & 部分 &
只标记''误差大''但\textbf{不解释为什么}；不发现误差的结构化模式 \\
MTP D-optimality & 人 (extrapolation grade) & 线性代数 & 部分 &
extrapolation grade 是几何度量，不关联误差的具体来源 \\
\textbf{SCX} & \textbf{算法 (误差分布)} & \textbf{自动筛选} &
\textbf{✅} & \textbf{误差驱动 + 可解释描述符 + 自动状态发现} \\
\end{longtable}

\paragraph{1.2.1 DIRECT sampling}\label{direct-sampling}

DIRECT (D-optimality-guided Iterative REgion Cluster Training) 是最接近
SCX 的方法：它对 SOAP 描述符做 PCA
降维后聚类，在每个簇中采样。但其关键限制：

\begin{itemize}
\tightlist
\item
  SOAP 描述符编码的是\textbf{局域化学环境的相似性}，而非误差模式
\item
  聚类边界由 SOAP 距离决定，而非模型失败边界
\item
  即使两个结构在 SOAP
  空间中相邻，它们的误差行为可能完全不同（例如：一个完美表面和一个部分重构的表面）
\end{itemize}

\paragraph{1.2.2 DP-GEN}\label{dp-gen}

DP-GEN 用 model deviation
作为不确定性度量，选择高偏差的构型加入训练集。其优点在于直接瞄准误差，但：

\begin{itemize}
\tightlist
\item
  \textbf{不解释}：\(\epsilon_{\text{dev}}\)
  是一个标量，告诉你''这个结构误差大''，但不告诉你''因为配位数小于 3.5
  且局域应变大于 3\%''
\item
  \textbf{无结构化状态}：DP-GEN
  产生的是一组''高偏差构型''，而非可解释的状态划分
\item
  这使得 DP-GEN
  的主动采样策略难以迁移到新材料------你只知道''更多数据更好''，但不知道''什么样的数据更好''
\end{itemize}

\paragraph{1.2.3 MTP D-optimality}\label{mtp-d-optimality}

MTP 用 D-optimality 准则衡量 extrapolation grade，但其：

\begin{itemize}
\tightlist
\item
  extrapolation grade
  是线性代数度量（设计矩阵的行列式），而非模型失败行为的直接度量
\item
  高 extrapolation grade
  并不意味着高误差------某些线性外推区域可能模型表现良好
\item
  同样不提供失败模式的物理解释
\end{itemize}

\paragraph{1.2.4
共同的根本缺陷}\label{ux5171ux540cux7684ux6839ux672cux7f3aux9677}

所有竞争方法都隐含地假设''人类知道什么是正确的状态划分''------但人类不知道。

\begin{itemize}
\tightlist
\item
  人知道 AlN 有 bulk、surface、defect 三种结构
\item
  人不知道的是：模型在''配位数 3.0-3.5 且局域应变 \textgreater{}
  3\%``的区域误差比平均值高 8 倍
\item
  后者才是真正需要成为独立状态的划分
\end{itemize}

\begin{center}\rule{0.5\linewidth}{0.5pt}\end{center}

\subsection{2.
两层描述符体系}\label{ux4e24ux5c42ux63cfux8ff0ux7b26ux4f53ux7cfb}

\subsubsection{2.1 设计原则}\label{ux8bbeux8ba1ux539fux5219}

两层描述符体系的设计源于一个简单但深刻的观察：

\begin{quote}
\textbf{失败的根源结构化了------但不是按人类熟悉的类别结构化的。}
\end{quote}

SCX
的两层描述符体系通过分离''人类知道的''和''算法发现的''，解决了这一问题。

\subsubsection{2.2 Layer
1：人类提供的粗粒度知识}\label{layer-1ux4ebaux7c7bux63d0ux4f9bux7684ux7c97ux7c92ux5ea6ux77e5ux8bc6}

第一层描述符是研究人员凭领域知识提供的材料''身份标签''，用于定义\textbf{初始采样空间}。

\begin{longtable}[]{@{}
  >{\raggedright\arraybackslash}p{(\linewidth - 6\tabcolsep) * \real{0.2973}}
  >{\raggedright\arraybackslash}p{(\linewidth - 6\tabcolsep) * \real{0.2432}}
  >{\raggedright\arraybackslash}p{(\linewidth - 6\tabcolsep) * \real{0.2432}}
  >{\raggedright\arraybackslash}p{(\linewidth - 6\tabcolsep) * \real{0.2162}}@{}}
\toprule\noalign{}
\begin{minipage}[b]{\linewidth}\raggedright
描述符类别
\end{minipage} & \begin{minipage}[b]{\linewidth}\raggedright
具体特征
\end{minipage} & \begin{minipage}[b]{\linewidth}\raggedright
物理意义
\end{minipage} & \begin{minipage}[b]{\linewidth}\raggedright
示例值
\end{minipage} \\
\midrule\noalign{}
\endhead
\bottomrule\noalign{}
\endlastfoot
化学组分 & 元素种类、化学计量比 & 化学空间定位 & AlN (1:1),
Al₀.₅Ga₀.₅N \\
晶体结构 & 空间群、晶格常数 & 热力学相 & wurtzite, zincblende,
rocksalt \\
热力学条件 & 温度、压力 & 相图坐标 & 300K-1800K, 0-20 GPa \\
形态类型 & 结构类别 & 材料形态 & bulk, surface, defect, interface,
cluster \\
计算参数 & k-points, cutoff & 计算设置 & 4x4x4, 500 eV \\
\end{longtable}

Layer 1 的关键函数在 \texttt{src/scx/encoders/base.py} 中定义为
\texttt{SCXStateEncoder} 抽象基类，其 \texttt{encode()} 和
\texttt{batch\_encode()} 方法提供统一接口。在
\texttt{src/scx/encoders/error\_driven.py} 第 98-114
行，\texttt{ErrorDrivenEncoder.encode()} 和 \texttt{batch\_encode()}
直接委托给 Layer 1 编码器：

\begin{Shaded}
\begin{Highlighting}[]
\KeywordTok{def}\NormalTok{ encode(}\VariableTok{self}\NormalTok{, x: Any) }\OperatorTok{{-}\textgreater{}}\NormalTok{ np.ndarray:}
    \ControlFlowTok{return} \VariableTok{self}\NormalTok{.layer1.encode(x)}

\KeywordTok{def}\NormalTok{ batch\_encode(}\VariableTok{self}\NormalTok{, X: }\BuiltInTok{list}\NormalTok{[Any]) }\OperatorTok{{-}\textgreater{}}\NormalTok{ np.ndarray:}
    \ControlFlowTok{return} \VariableTok{self}\NormalTok{.layer1.batch\_encode(X)}
\end{Highlighting}
\end{Shaded}

\textbf{Layer 1
的职责}：不是定义最终的状态划分，而是提供一个\textbf{候选描述符池}------后续
Layer 2 从池中选择与误差相关的子集。

\subsubsection{2.3 Layer
2：算法发现的错误驱动状态}\label{layer-2ux7b97ux6cd5ux53d1ux73b0ux7684ux9519ux8befux9a71ux52a8ux72b6ux6001}

第二层描述符不是在研究开始前预设的，而是算法在误差分析阶段\textbf{自动计算、筛选、组合}的。

\begin{longtable}[]{@{}
  >{\raggedright\arraybackslash}p{(\linewidth - 4\tabcolsep) * \real{0.3333}}
  >{\raggedright\arraybackslash}p{(\linewidth - 4\tabcolsep) * \real{0.2727}}
  >{\raggedright\arraybackslash}p{(\linewidth - 4\tabcolsep) * \real{0.3939}}@{}}
\toprule\noalign{}
\begin{minipage}[b]{\linewidth}\raggedright
描述符类别
\end{minipage} & \begin{minipage}[b]{\linewidth}\raggedright
候选特征
\end{minipage} & \begin{minipage}[b]{\linewidth}\raggedright
误差来源假设
\end{minipage} \\
\midrule\noalign{}
\endhead
\bottomrule\noalign{}
\endlastfoot
配位数 (CN) & CN@r=2.5, CN@r=3.0, CN@r=4.0 Å &
表面/低配位区域力场不准 \\
键拉伸 & max bond stretch, mean bond stretch & 远离平衡键长的能量误差 \\
局域应变 & 应变张量不变量 J₁, J₂ & 高应变下的弹性非线性 \\
杂化比例 & sp²/sp³ 比率, 键角分布矩 & 杂化转变区域的势能面交叉 \\
位移幅度 & 原子位移范数 & 高温热扰动的非谐效应 \\
配位变化 & Voronoi 体积, 配位多面体畸变 & 缺陷/相变前驱体 \\
ACE 内部分量 & ACE basis 各分量贡献 & 特定基函数的系统误差 \\
力的角分布 & 切向/径向力误差比 & 各向异性环境中的力场偏差 \\
\end{longtable}

Layer 2 的核心逻辑在 \texttt{src/scx/encoders/error\_driven.py} 的
\texttt{ErrorDrivenEncoder} 类（第 38-492 行）。其方法
\texttt{fit\_error\_states()}（第 149-221
行）实现了完整的误差驱动状态发现流程。

\subsubsection{2.4 两层的关系}\label{ux4e24ux5c42ux7684ux5173ux7cfb}

两层之间不是并列关系，而是\textbf{层级选择和细化的关系}：

\begin{verbatim}
研究人员："我要做 AlGaN，包括 bulk/surface/defect，温度范围 300-1800K"
         |
         v
Layer 1：生成初始 DFT 采样网格（按人的定义分类）
         |
         v
训练 Single ACE baseline
         |
         v
Layer 2：算法计算 per-atom error -> 映射到候选描述符空间
         自动筛选最相关描述符 -> 在误差相关子空间中聚类
         |
         v
算法发现："在 CN<3.5 且 local strain>3% 的区域，力误差是平均水平的 8 倍"
         |
         v
算法建议："生成 CN=2.5-3.5, local strain=2-5% 的结构，约 50 个 DFT 计算"
         |
         v
研究人员跑 DFT -> 加入训练 -> 误差下降
         |
         v
迭代直到误差景观平坦
\end{verbatim}

\textbf{用 SCX 状态空间的术语说}：Layer 1 定义初始状态空间
\(\mathcal{S}_1\)，Layer 2 在每个 \(s \in \mathcal{S}_1\) 中发现子状态
\(\mathcal{S}_2(s)\)，最终的细粒度状态空间为：

\[\mathcal{S}_{\text{final}} = \bigcup_{s \in \mathcal{S}_1} \mathcal{S}_2(s)\]

代码层面，\texttt{TwoLayerStateDiscovery}
类（\texttt{src/scx/state/two\_layer.py} 第 45-451
行）实现了这一级联。其 \texttt{discover()} 方法（第 77-187 行）先做
Layer 1 聚类，再调用 \texttt{ErrorDrivenEncoder.fit\_error\_states()} 做
Layer 2 的误差驱动细化。

\begin{center}\rule{0.5\linewidth}{0.5pt}\end{center}

\subsection{3. 数学形式化}\label{ux6570ux5b66ux5f62ux5f0fux5316}

\subsubsection{3.1 符号体系}\label{ux7b26ux53f7ux4f53ux7cfb}

\begin{longtable}[]{@{}
  >{\raggedright\arraybackslash}p{(\linewidth - 4\tabcolsep) * \real{0.3333}}
  >{\raggedright\arraybackslash}p{(\linewidth - 4\tabcolsep) * \real{0.3333}}
  >{\raggedright\arraybackslash}p{(\linewidth - 4\tabcolsep) * \real{0.3333}}@{}}
\toprule\noalign{}
\begin{minipage}[b]{\linewidth}\raggedright
符号
\end{minipage} & \begin{minipage}[b]{\linewidth}\raggedright
含义
\end{minipage} & \begin{minipage}[b]{\linewidth}\raggedright
维度
\end{minipage} \\
\midrule\noalign{}
\endhead
\bottomrule\noalign{}
\endlastfoot
\(\varphi_1(x)\) & Layer 1 编码 & \(\mathbb{R}^d\) \\
\(r(x)\) & 残差（误差） & \(\mathbb{R}\) 或 \(\mathbb{R}^M\) \\
\(\text{MI}(\varphi_1[:, j], r)\) & 第 j 维与残差的互信息 &
\(\mathbb{R}\) \\
\(D^*\) & 误差相关维度索引集 & \(\{1,\dots,d\}\) 的子集 \\
\(\tau\) & 互信息选择阈值 & \(\mathbb{R}^+\) \\
\(K\) & 目标状态数 & \(\mathbb{N}\) \\
\(\mathcal{S}_{\text{human}}\) & 人类定义的状态划分 & 从 \(\mathcal{X}\)
到 \(\{1,\dots,K\}\) 的映射 \\
\(\mathcal{S}^*\) & 误差驱动发现的状态划分 & 同上 \\
\end{longtable}

\subsubsection{3.2 Error-Relevant
Subspace}\label{error-relevant-subspace}

\textbf{定义 3.1（误差相关子空间）}。给定 Layer 1 编码矩阵
\(\Phi_1 \in \mathbb{R}^{N \times d}\) 和残差向量
\(r \in \mathbb{R}^N\)，误差相关子空间定义为与残差互信息超过阈值的维度索引集：

\[D^* = \{j: \text{MI}(\Phi_1[:, j], r) > \tau\}\]

其中互信息定义为：

\[\text{MI}(X_j, r) = \int \int p(x_j, r) \log \frac{p(x_j, r)}{p(x_j)p(r)} \, dx_j \, dr\]

在实现中（\texttt{src/scx/encoders/error\_driven.py} 第 248-270
行），互信息通过
\texttt{sklearn.feature\_selection.mutual\_info\_regression}
估计，并归一化到 \([0,1]\)：

\begin{Shaded}
\begin{Highlighting}[]
\NormalTok{mi }\OperatorTok{=}\NormalTok{ mutual\_info\_regression(phi, residuals, random\_state}\OperatorTok{=}\DecValTok{42}\NormalTok{)}
\NormalTok{mi }\OperatorTok{=}\NormalTok{ np.nan\_to\_num(mi, nan}\OperatorTok{=}\FloatTok{0.0}\NormalTok{, posinf}\OperatorTok{=}\FloatTok{0.0}\NormalTok{, neginf}\OperatorTok{=}\FloatTok{0.0}\NormalTok{)}
\NormalTok{mi }\OperatorTok{=}\NormalTok{ mi }\OperatorTok{/}\NormalTok{ mi.}\BuiltInTok{max}\NormalTok{()  }\CommentTok{\# normalize to [0,1]}
\end{Highlighting}
\end{Shaded}

\textbf{选择机制}（第 315-326 行）：选择互信息最大的 top-\(k\) 维度：

\begin{Shaded}
\begin{Highlighting}[]
\NormalTok{n\_select }\OperatorTok{=} \BuiltInTok{max}\NormalTok{(}\DecValTok{3}\NormalTok{, }\BuiltInTok{int}\NormalTok{(d }\OperatorTok{*} \VariableTok{self}\NormalTok{.selected\_dims\_ratio))}
\NormalTok{top\_dims }\OperatorTok{=}\NormalTok{ np.argsort(importance)[}\OperatorTok{{-}}\NormalTok{n\_select:]}
\end{Highlighting}
\end{Shaded}

默认 \texttt{selected\_dims\_ratio\ =\ 1/3}，意味着选择与误差最相关的前
1/3 维度。

\textbf{定理 3.1（误差相关子空间的最优性）}。在所有与 \(\Phi_1\)
同维度的子空间选择方案中，\(D^* = \{j: \text{MI}(\Phi_1[:, j], r) > \tau\}\)
最大化所选维度对残差 \(r\) 的总信息。

\textbf{证明}。对于任何维度子集 \(S \subseteq \{1,\dots,d\}\)，总信息由
\(\sum_{j \in S} \text{MI}(\Phi_1[:, j], r)\)
给出（假设维度间独立，或在弱相关假设下作为近似）。选择 \(\text{MI}\)
值最大的 \(k\) 个维度等价于解：

\[D^*_k = \arg\max_{S: |S|=k} \sum_{j \in S} \text{MI}(\Phi_1[:, j], r)\]

这是子集选择的贪心最优解（对非负独立项，排序选择即全局最优）。\(\square\)

\subsubsection{3.3 Error-Driven States}\label{error-driven-states}

\textbf{定义 3.2（误差驱动状态）}。在误差相关子空间 \(D^*\) 上的
\(K\)-聚类即为误差驱动状态：

\[\mathcal{S}^* = \text{Cluster}(\Phi_1[:, D^*], K)\]

其中 \(\text{Cluster}(\cdot, K)\)
为任意聚类算法（K-means、HDBSCAN、谱聚类）。在 SCX
实现中（\texttt{src/scx/encoders/error\_driven.py} 第 188-190
行），聚类委托给 Layer 1 编码器：

\begin{Shaded}
\begin{Highlighting}[]
\NormalTok{phi\_error }\OperatorTok{=}\NormalTok{ phi[:, top\_dims] }\ControlFlowTok{if} \BuiltInTok{len}\NormalTok{(top\_dims) }\OperatorTok{\textgreater{}} \DecValTok{0} \ControlFlowTok{else}\NormalTok{ phi}
\NormalTok{n\_clusters }\OperatorTok{=} \BuiltInTok{min}\NormalTok{(}\VariableTok{self}\NormalTok{.n\_error\_states, }\BuiltInTok{max}\NormalTok{(}\DecValTok{1}\NormalTok{, n\_samples }\OperatorTok{{-}} \DecValTok{1}\NormalTok{))}
\NormalTok{labels, centroids }\OperatorTok{=} \VariableTok{self}\NormalTok{.layer1.cluster(phi\_error, n\_clusters)}
\end{Highlighting}
\end{Shaded}

\textbf{关键约束}：聚类空间 \(\Phi_1[:, D^*]\) 的维度远小于原始空间
\(\Phi_1\)（默认选择 1/3 维度），这意味着：

\[\dim(\Phi_1[:, D^*]) \ll \dim(\Phi_1)\]

这不仅是降维，更是\textbf{语义过滤}------排除与误差无关的结构信息，保留与误差相关的失败模式。

\subsubsection{3.4 优势定理}\label{ux4f18ux52bfux5b9aux7406}

\textbf{定理 3.2（误差驱动状态占优人类定义状态）}。令
\(\mathcal{S}_{\text{human}}\)
为人类先验知识定义的状态划分（如按晶体结构、化学组分分类），\(\mathcal{S}^*\)
为误差驱动发现的状态划分。在相同的状态数 \(K = |\mathcal{S}|\) 下：

\[\mathbb{E}_{s \sim \mathcal{S}^*}[\text{Var}(r \mid s)] \leq \mathbb{E}_{s \sim \mathcal{S}_{\text{human}}}[\text{Var}(r \mid s)]\]

即 error-driven states 的组内残差方差不超过人类定义状态的组内残差方差。

\textbf{证明}。

\textbf{引理 3.2.1}。对人类定义的任意状态
\(s \in \mathcal{S}_{\text{human}}\)，其在 \(\varphi_1\) 空间中的表示
\(\mu_s = \mathbb{E}[\varphi_1(x) \mid x \in s]\) 可能包含大量与 \(r\)
无关的维度。

\textbf{证明}。人类定义状态时使用的特征
\(\psi(x)\)（如晶体结构、化学组分）可能与 \(\varphi_1\) 存在非线性关系
\(\varphi_1(x) = g(\psi(x)) + \epsilon\)。如果 \(g\) 的某些分量与 \(r\)
无关，则这些分量在聚类中作为噪声存在，增加组内误差方差。

\textbf{引理 3.2.2}。K-means 在投影子空间 \(D^*\) 上的目标函数：

\[\min_{\mathcal{S}} \sum_{s \in \mathcal{S}} \sum_{x \in s} \|\varphi_1(x)[D^*] - \mu_s\|^2\]

等价于在残差相关维度上的方差最小化。其中
\(\mu_s = \frac{1}{|s|} \sum_{x \in s} \varphi_1(x)[D^*]\)。

\textbf{证明}。由于 \(D^*\) 的每个维度 \(j\) 满足
\(\text{MI}(\varphi_1[:, j], r) > \tau\)，在 \(\varphi_1[:, D^*]\)
上的聚类直接捕捉了与 \(r\) 最相关的结构。K-means
的组内平方和最小化等价于在每个簇内最小化特征方差，而由于这些特征与 \(r\)
相关，特征方差的减小意味着 \(r\) 条件方差的减小。

\textbf{主定理证明}。由 \(\mathcal{S}^*\) 的定义，\(\mathcal{S}^*\)
直接优化 \(\text{Var}(r \mid s)\) 的上界（通过在 \(D^*\)
中聚类）。\(\mathcal{S}_{\text{human}}\) 在 \(\varphi_1\)
的全空间上聚类，包含了与 \(r\) 无关的''噪声维度''。由引理 3.2.1 和
3.2.2，\(\mathcal{S}^*\) 在 \(r\) 上的条件方差不大于
\(\mathcal{S}_{\text{human}}\)
的条件方差。等号成立当且仅当人类定义的状态恰好在 \(D^*\)
中是可分的。\(\square\)

\textbf{推论 3.2.1（严格优势）}。如果 \(\mathcal{S}_{\text{human}}\)
不是基于 \(r\) 的充分统计量，即存在
\(u \neq v \in \mathcal{S}_{\text{human}}\) 使得
\(\mathbb{E}[r \mid u] \neq \mathbb{E}[r \mid v]\) 但
\(\mu_u[D^*] = \mu_v[D^*]\)的概率小于 1，则不等号严格成立：

\[\mathbb{E}_{s \sim \mathcal{S}^*}[\text{Var}(r \mid s)] < \mathbb{E}_{s \sim \mathcal{S}_{\text{human}}}[\text{Var}(r \mid s)]\]

\subsubsection{3.5
与信息瓶颈的关系}\label{ux4e0eux4fe1ux606fux74f6ux9888ux7684ux5173ux7cfb}

两层状态发现可以理解为\textbf{条件信息瓶颈}的迭代求解。信息瓶颈原理
(Tishby et al., 1999) 给出了最优表示 \(\tilde{X}\) 的变分目标：

\[\min_{\Pi} I(X; \tilde{X}) - \beta I(\tilde{X}; Y)\]

在 SCX 的两层框架中，这一原理被条件化为：

\[\min_{\Pi} I(\varphi_1; \tilde{X} \mid \mathcal{F}) - \beta I(\tilde{X}; r \mid \mathcal{F})\]

其中 \(\mathcal{F} = \{f_1, \dots, f_M\}\) 为专家集合，\(r\)
为残差向量。与标准信息瓶颈的关键区别：

\begin{enumerate}
\def\labelenumi{\arabic{enumi}.}
\tightlist
\item
  \textbf{以 \(\varphi_1\) 替代 \(X\)}：表示从原始输入空间提升到 Layer 1
  嵌入空间，降低复杂度
\item
  \textbf{以残差 \(r\) 替代标签
  \(Y\)}：不是预测''结构是什么''，而是预测''模型在哪里失败''
\item
  \textbf{以专家 \(\mathcal{F}\)
  为条件}：状态划分依赖于当前的专家集合，当专家更新时状态可演化
\end{enumerate}

条件信息瓶颈的迭代求解通过 SCX 的 Phase A-E
流程序列实现。每轮迭代等价于一次信息瓶颈的变分更新： - Phase B 的聚类
\(\approx\) 压缩 \(I(\varphi_1; \tilde{X} \mid \mathcal{F})\) - Phase C
的定向采样 \(\approx\) 增加 \(I(\tilde{X}; r \mid \mathcal{F})\) - Phase
D 的模型更新 \(\approx\) 更新 \(\mathcal{F}\)，改变条件分布

\subsubsection{3.6
状态演化的动力学}\label{ux72b6ux6001ux6f14ux5316ux7684ux52a8ux529bux5b66}

这是两层描述符框架独有的理论特征------状态划分不固定。

\textbf{定义 3.3（状态演化算子）}。令 \(\mathcal{F}^{(t)}\) 为第 \(t\)
轮迭代后的专家集合，\(\tilde{X}^{(t)}\) 为对应的状态划分。状态演化算子
\(T\) 定义为：

\[\tilde{X}^{(t+1)} = T(\tilde{X}^{(t)} \mid \mathcal{F}^{(t+1)}, r^{(t+1)})\]

其中 \(r^{(t+1)}\) 是更新后的模型误差。

\textbf{学习轨迹}：

\begin{verbatim}
t=0: 人类定义状态 S_human ≈ 信息瓶颈的"先验"状态
t=1: 在 baseline 误差上聚类 → S*(1)，发现"CN<3.5 区域误差高"
t=2: 在该区域补充 DFT → 误差下降 → 误差景观变化
t=3: 重新聚类 → S*(2)，可能发现"新的高误差区域：sp2 过渡态"
...
t=∞: 误差景观平坦 → 状态划分稳定（收敛）
\end{verbatim}

\textbf{定理 3.3（状态收敛性）}。如果误差景观的 Lipschitz 常数 \(L_{r}\)
有限，且每次迭代补充的 DFT 样本数 \(N_t\) 满足
\(\sum_{t=1}^\infty N_t = \infty\)，则状态划分 \(\tilde{X}^{(t)}\)
几乎必然收敛到一个固定点。

\textbf{证明概要}。将状态演化视为 \(\varphi_1\) 空间中一个 Markov
链。由于每轮迭代补充的 DFT 数据覆盖了高误差区域，\(r^{(t)}\)
逐点递减。当 \(r^{(t)}\) 收敛到平稳过程时，\(D^*\)
和聚类结果随之稳定。\(\square\)

\begin{center}\rule{0.5\linewidth}{0.5pt}\end{center}

\subsection{4. SCX 实现}\label{scx-ux5b9eux73b0}

\subsubsection{4.1 ErrorDrivenEncoder}\label{errordrivenencoder}

文件：\texttt{src/scx/encoders/error\_driven.py}，第 38-492 行。

这是两层描述符体系的核心实现，负责从误差分布中自动发现状态。

\paragraph{4.1.1 初始化（第 66-92
行）}\label{ux521dux59cbux5316ux7b2c-66-92-ux884c}

\begin{Shaded}
\begin{Highlighting}[]
\KeywordTok{class}\NormalTok{ ErrorDrivenEncoder(SCXStateEncoder):}
    \KeywordTok{def} \FunctionTok{\_\_init\_\_}\NormalTok{(}
        \VariableTok{self}\NormalTok{,}
\NormalTok{        layer1\_encoder: SCXStateEncoder,}
\NormalTok{        n\_error\_states: }\BuiltInTok{int} \OperatorTok{=} \DecValTok{10}\NormalTok{,}
\NormalTok{        feature\_selection: }\BuiltInTok{str} \OperatorTok{=} \StringTok{"mutual\_info"}\NormalTok{,}
\NormalTok{        selected\_dims\_ratio: }\BuiltInTok{float} \OperatorTok{=} \FloatTok{1.0} \OperatorTok{/} \FloatTok{3.0}\NormalTok{,}
\NormalTok{    ) }\OperatorTok{{-}\textgreater{}} \VariableTok{None}\NormalTok{:}
\end{Highlighting}
\end{Shaded}

关键参数： - \texttt{layer1\_encoder}：第一层编码器，提供候选描述符池 -
\texttt{n\_error\_states}：目标状态数（可自动调整） -
\texttt{feature\_selection}：特征选择方法，支持
\texttt{mutual\_info}（默认）、\texttt{correlation}、\texttt{shap}、\texttt{none}
- \texttt{selected\_dims\_ratio}：保留与误差最相关的维度比例，默认 1/3

\paragraph{4.1.2 核心方法：fit\_error\_states（第 149-221
行）}\label{ux6838ux5fc3ux65b9ux6cd5fit_error_statesux7b2c-149-221-ux884c}

完整流程：

\begin{enumerate}
\def\labelenumi{\arabic{enumi}.}
\tightlist
\item
  \textbf{Layer 1 编码}（第 176
  行）：\texttt{self.\_layer1\_phi\_\ =\ self.batch\_encode(X\_raw)}
  将原始输入映射到 \(\varphi_1\) 空间
\item
  \textbf{特征重要性计算}（第 180
  行）：\texttt{feature\_importance\ =\ self.\_compute\_feature\_importance(phi,\ residuals)}
  计算每个维度与残差的相关性
\item
  \textbf{维度选择}（第 183-186
  行）：\texttt{top\_dims\ =\ self.\_select\_top\_dims(feature\_importance,\ phi.shape{[}1{]})}
  保留与误差最相关的维度
\item
  \textbf{误差相关子空间聚类}（第 189-190 行）：在 \(\Phi_1[:, D^*]\)
  上聚类，发现 error states
\item
  \textbf{Layer 1 细化}（第 194-197 行）：如果提供了
  \texttt{layer1\_labels}，在每个 Layer 1 簇内部进一步拆分子状态
\item
  \textbf{状态描述生成}（第 203-217 行）：为每个 error state
  计算特征轮廓、平均误差、样本数和自动描述
\end{enumerate}

\paragraph{4.1.3
特征选择方法}\label{ux7279ux5f81ux9009ux62e9ux65b9ux6cd5}

\begin{longtable}[]{@{}
  >{\raggedright\arraybackslash}p{(\linewidth - 6\tabcolsep) * \real{0.2000}}
  >{\raggedright\arraybackslash}p{(\linewidth - 6\tabcolsep) * \real{0.2000}}
  >{\raggedright\arraybackslash}p{(\linewidth - 6\tabcolsep) * \real{0.3000}}
  >{\raggedright\arraybackslash}p{(\linewidth - 6\tabcolsep) * \real{0.3000}}@{}}
\toprule\noalign{}
\begin{minipage}[b]{\linewidth}\raggedright
方法
\end{minipage} & \begin{minipage}[b]{\linewidth}\raggedright
度量
\end{minipage} & \begin{minipage}[b]{\linewidth}\raggedright
计算方式
\end{minipage} & \begin{minipage}[b]{\linewidth}\raggedright
适用场景
\end{minipage} \\
\midrule\noalign{}
\endhead
\bottomrule\noalign{}
\endlastfoot
\texttt{mutual\_info} & 互信息 &
\texttt{sklearn.feature\_selection.mutual\_info\_regression} &
默认，捕捉非线性关系 \\
\texttt{correlation} & Pearson 相关系数 & \(|\text{corr}(\phi_j, r)|\) &
快速近似，仅线性关系 \\
\texttt{shap} & SHAP 值 & \texttt{shap.TreeExplainer(RandomForest)} &
模型感知的特征归因 \\
\texttt{none} & 均匀权重 & 所有维度保留 & 对比实验，不筛选 \\
\end{longtable}

第 240-243 行的路由逻辑：

\begin{Shaded}
\begin{Highlighting}[]
\ControlFlowTok{if} \VariableTok{self}\NormalTok{.feature\_selection }\OperatorTok{==} \StringTok{"mutual\_info"}\NormalTok{:}
    \ControlFlowTok{return} \VariableTok{self}\NormalTok{.\_mutual\_info\_regression(phi, residuals)}
\ControlFlowTok{elif} \VariableTok{self}\NormalTok{.feature\_selection }\OperatorTok{==} \StringTok{"correlation"}\NormalTok{:}
    \ControlFlowTok{return} \VariableTok{self}\NormalTok{.\_correlation\_importance(phi, residuals)}
\ControlFlowTok{elif} \VariableTok{self}\NormalTok{.feature\_selection }\OperatorTok{==} \StringTok{"shap"}\NormalTok{:}
    \ControlFlowTok{return} \VariableTok{self}\NormalTok{.\_shap\_importance(phi, residuals)}
\end{Highlighting}
\end{Shaded}

\textbf{互信息的选择理由}：与 Pearson
相关系数不同，互信息可以捕捉非线性关系（这是必须的------\(r\) 与
\(\varphi_1\) 之间的关系可能是高度非线性的）。与 SHAP
相比，互信息无需训练代理模型，计算更快，且对 \(\varphi_1\)
的分布形式无假设。

\paragraph{4.1.4 高误差状态识别（第 419-446
行）}\label{ux9ad8ux8befux5deeux72b6ux6001ux8bc6ux522bux7b2c-419-446-ux884c}

\texttt{get\_high\_error\_states()} 返回平均误差超过阈值（默认全局 mean
+ 1 std）的状态 ID。这对定向采样决策至关重要：

\begin{Shaded}
\begin{Highlighting}[]
\KeywordTok{def}\NormalTok{ get\_high\_error\_states(}\VariableTok{self}\NormalTok{, threshold}\OperatorTok{=}\VariableTok{None}\NormalTok{):}
\NormalTok{    errors }\OperatorTok{=}\NormalTok{ np.array([s[}\StringTok{"mean\_error"}\NormalTok{] }\ControlFlowTok{for}\NormalTok{ s }\KeywordTok{in} \VariableTok{self}\NormalTok{.error\_states\_.values()])}
    \ControlFlowTok{if}\NormalTok{ threshold }\KeywordTok{is} \VariableTok{None}\NormalTok{:}
\NormalTok{        threshold }\OperatorTok{=}\NormalTok{ errors.mean() }\OperatorTok{+}\NormalTok{ errors.std()}
    \ControlFlowTok{return}\NormalTok{ [sid }\ControlFlowTok{for}\NormalTok{ sid, s }\KeywordTok{in} \VariableTok{self}\NormalTok{.error\_states\_.items()}
            \ControlFlowTok{if}\NormalTok{ s[}\StringTok{"mean\_error"}\NormalTok{] }\OperatorTok{\textgreater{}}\NormalTok{ threshold]}
\end{Highlighting}
\end{Shaded}

\subsubsection{4.2 TwoLayerStateDiscovery}\label{twolayerstatediscovery}

文件：\texttt{src/scx/state/two\_layer.py}，第 45-451 行。

\paragraph{4.2.1 完整的两层管线（第 77-187
行）}\label{ux5b8cux6574ux7684ux4e24ux5c42ux7ba1ux7ebfux7b2c-77-187-ux884c}

\texttt{discover()} 方法将 Layer 1 和 Layer 2 整合为一个完整的管线：

\begin{verbatim}
    输入: X_raw, residuals
         |
         v
    Layer 1: batch_encode(X_raw) -> phi
    StateDiscovery.cluster(phi) -> layer1_labels
         |
         v
    构建 layer1 状态摘要
    (n_samples, proportion, mean_residual, centroid)
         |
         v
    Layer 2: ErrorDrivenEncoder.fit_error_states(
        X_raw, residuals, layer1_labels)
         |
         v
    生成定向采样建议
    recommend_sampling(error_states, budget)
         |
         v
    输出: layer1_labels, layer2_labels, error_states,
          layer1_states, recommendations, metrics
\end{verbatim}

\paragraph{4.2.2 定向采样建议（第 193-308
行）}\label{ux5b9aux5411ux91c7ux6837ux5efaux8baeux7b2c-193-308-ux884c}

\texttt{recommend\_sampling()} 基于 error state
分析生成结构补充建议。优先级逻辑（第 247-260 行）：

\begin{itemize}
\tightlist
\item
  \textbf{高误差 + 高密度} = 高优先级（快速改善模型）
\item
  \textbf{高误差 + 低密度} = 中优先级（可能是稀疏区域）
\item
  \textbf{低误差} = 低优先级
\item
  \textbf{极高误差 + 极少样本} = 标记为噪声检查
\end{itemize}

优先级分数（第 251-252 行）：

\[\text{priority}(s) = \bar{r}_s \cdot (1 + \sqrt{N_s / N_{\text{total}}})\]

其中 \(\bar{r}_s\) 为状态 \(s\) 的平均误差，\(N_s\) 为样本数。密度因子
\(\sqrt{N_s / N_{\text{total}}}\)
确保高密度区域获得更多采样（因为高密度区域改善模型的效果更显著）。

高误差状态的优先级加倍（第 254-255 行）：

\begin{Shaded}
\begin{Highlighting}[]
\ControlFlowTok{if}\NormalTok{ e\_s }\OperatorTok{\textgreater{}}\NormalTok{ err\_mean }\OperatorTok{+} \FloatTok{1.5} \OperatorTok{*}\NormalTok{ err\_std }\KeywordTok{and}\NormalTok{ n\_s }\OperatorTok{\textgreater{}=}\NormalTok{ min\_samples\_per\_state:}
\NormalTok{    priority }\OperatorTok{*=} \FloatTok{2.0}
\end{Highlighting}
\end{Shaded}

\paragraph{4.2.3 对比分析（第 314-416
行）}\label{ux5bf9ux6bd4ux5206ux6790ux7b2c-314-416-ux884c}

\texttt{compare\_with\_pure\_layer1()} 提供严格的对照实验：

\begin{itemize}
\tightlist
\item
  纯 Layer 1 聚类（在 \(\varphi_1\) 全空间聚类）
\item
  两层聚类（Layer 1 + ErrorDrivenEncoder）
\item
  对比指标：\texttt{state\_error\_spread}（各状态平均误差的方差）、\texttt{max\_min\_error\_ratio}（最大/最小状态误差比）
\end{itemize}

\textbf{预期结果}：两层的 \texttt{error\_spread} 和
\texttt{max\_min\_error\_ratio} 应显著低于纯 Layer
1，证明误差驱动状态划分使得组内误差更均匀。

\subsubsection{4.3 与 SCX Core
的架构映射}\label{ux4e0e-scx-core-ux7684ux67b6ux6784ux6620ux5c04}

两层描述符体系与 SCX 核心框架的模块对应关系：

\begin{longtable}[]{@{}
  >{\raggedright\arraybackslash}p{(\linewidth - 4\tabcolsep) * \real{0.3111}}
  >{\raggedright\arraybackslash}p{(\linewidth - 4\tabcolsep) * \real{0.4889}}
  >{\raggedright\arraybackslash}p{(\linewidth - 4\tabcolsep) * \real{0.2000}}@{}}
\toprule\noalign{}
\begin{minipage}[b]{\linewidth}\raggedright
SCX Core 模块
\end{minipage} & \begin{minipage}[b]{\linewidth}\raggedright
两层描述符体系中的角色
\end{minipage} & \begin{minipage}[b]{\linewidth}\raggedright
对应代码
\end{minipage} \\
\midrule\noalign{}
\endhead
\bottomrule\noalign{}
\endlastfoot
\texttt{StateEncoder} & Layer 1 编码器 &
\texttt{scx.encoders.base.SCXStateEncoder} \\
\texttt{StateDiscovery} & Layer 1 初始状态划分 &
\texttt{scx.state.discovery.StateDiscovery} \\
\texttt{ErrorDrivenEncoder} & Layer 2 误差驱动编码器 &
\texttt{scx.encoders.error\_driven.ErrorDrivenEncoder} \\
\texttt{TwoLayerStateDiscovery} & 完整的两层管线 &
\texttt{scx.state.two\_layer.TwoLayerStateDiscovery} \\
\texttt{ExpertReliability} & 状态条件专家可靠性估计 &
利用两层状态划分计算 \(R_m(s)\) \\
\texttt{ExpertRouter} & 专家路由选择 & 使用 error-driven states
作为路由空间 \\
\texttt{AcquisitionStrategy} & 定向采样策略 & 继承
\texttt{recommend\_sampling()} 的输出 \\
\end{longtable}

在 \texttt{discover()} 的输出中，\texttt{layer2\_labels} 直接输入到
\texttt{ExpertReliability} 模块计算条件风险 \(R_m(s)\)，然后输入到
\texttt{ExpertRouter}
做专家路由。这确保了\textbf{路由空间与误差结构对齐}------这正是 SCX
框架的核心假设。

\subsubsection{4.4
当前实现的局限}\label{ux5f53ux524dux5b9eux73b0ux7684ux5c40ux9650}

\begin{enumerate}
\def\labelenumi{\arabic{enumi}.}
\item
  \textbf{聚类算法不可配置}：\texttt{ErrorDrivenEncoder} 的聚类硬编码为
  \texttt{self.layer1.cluster()}（第 190 行）。当 Layer 1 基于 K-means
  时，Layer 2 使用的是同一 K-means，但误差相关子空间可能需要密度聚类（如
  HDBSCAN）。
\item
  \textbf{互信息对离群值敏感}：\texttt{mutual\_info\_regression} 用 k-NN
  估计互信息，对样本量小或离群值多的情况不稳定（内部回退机制在第 268
  行用 \texttt{try/except} 处理）。
\item
  \textbf{Layer 2 细化不够灵活}：\texttt{\_refine\_with\_layer1()}
  方法（第 332-385 行）仅在 Layer 1
  簇内误差方差大时做二分拆分，不支持更细粒度的子结构发现。
\item
  \textbf{描述自动生成原始}：\texttt{\_describe\_state()}（第 392-413
  行）仅输出维度索引和均值，不提供物理解释。这对于论文写作或研究人员理解而言不够直观。
\end{enumerate}

\begin{center}\rule{0.5\linewidth}{0.5pt}\end{center}

\subsection{5.
与竞争方法的对比}\label{ux4e0eux7adeux4e89ux65b9ux6cd5ux7684ux5bf9ux6bd4}

\subsubsection{5.1 系统化对比}\label{ux7cfbux7edfux5316ux5bf9ux6bd4}

\begin{longtable}[]{@{}
  >{\raggedright\arraybackslash}p{(\linewidth - 8\tabcolsep) * \real{0.1000}}
  >{\raggedright\arraybackslash}p{(\linewidth - 8\tabcolsep) * \real{0.1333}}
  >{\raggedright\arraybackslash}p{(\linewidth - 8\tabcolsep) * \real{0.1333}}
  >{\raggedright\arraybackslash}p{(\linewidth - 8\tabcolsep) * \real{0.2833}}
  >{\raggedright\arraybackslash}p{(\linewidth - 8\tabcolsep) * \real{0.3500}}@{}}
\toprule\noalign{}
\begin{minipage}[b]{\linewidth}\raggedright
维度
\end{minipage} & \begin{minipage}[b]{\linewidth}\raggedright
DIRECT
\end{minipage} & \begin{minipage}[b]{\linewidth}\raggedright
DP-GEN
\end{minipage} & \begin{minipage}[b]{\linewidth}\raggedright
MTP D-optimality
\end{minipage} & \begin{minipage}[b]{\linewidth}\raggedright
\textbf{SCX + 两层描述符}
\end{minipage} \\
\midrule\noalign{}
\endhead
\bottomrule\noalign{}
\endlastfoot
\textbf{状态定义者} & 人 (SOAP 聚类) & 人 (model deviation 阈值) & 人
(extrapolation grade) & \textbf{算法 (误差分布聚类)} \\
\textbf{描述符来源} & 人类预设 SOAP & model deviation 标量 &
线性代数设计矩阵 & \textbf{算法自动筛选} \\
\textbf{误差驱动？} & ❌ & 部分 (标量标记) & 部分 (几何度量) &
\textbf{✅ 完全误差驱动} \\
\textbf{状态可解释性} & SOAP 成分解释 & 无 & 无 (extrapolation grade
无物理意义) & \textbf{互信息 + 特征轮廓 → 物理解释} \\
\textbf{状态演化} & 固定 (预设聚类) & 固定 (阈值) & 固定 (D-optimality)
& \textbf{动态 (随模型迭代)} \\
\textbf{采样策略} & 簇内随机 & 高偏差构型 & 高 extrapolation &
\textbf{误差引导定向采样} \\
\textbf{对 DF 的需求} & 低 (一次性聚类) & 中 (需 MD 轨迹) & 中
(需线性代数) & \textbf{中 (需误差计算)} \\
\textbf{专家路由} & 无 & 无 & 无 & \textbf{状态条件专家路由} \\
\textbf{最适合场景} & 初始化采样 & 迭代训练增强 & 线性基函数拟合 &
\textbf{异质化学环境 + 多目标} \\
\end{longtable}

\subsubsection{5.2 DIRECT
的局限性}\label{direct-ux7684ux5c40ux9650ux6027}

DIRECT 是最接近 SCX+两层描述符的竞争方法，但根本区别在于聚类空间的选取：

\begin{itemize}
\tightlist
\item
  DIRECT 用 SOAP-PCA 空间聚类 → 优化的是''结构相似性''
\item
  SCX 用 \(\Phi_1[:, D^*]\) 空间聚类 → 优化的是''误差模式相似性''
\end{itemize}

\textbf{两种空间的几何差异}：设 \(d_{\text{SOAP}}(x_i, x_j)\) 为 SOAP
距离，\(d_{\text{error}}(x_i, x_j) = \|\phi(x_i)[D^*] - \phi(x_j)[D^*]\|\)
为误差空间距离。存在结构相似但误差不相似的情况（如：两个不同表面的 SOAP
距离小但力误差差异大），也存在结构不相似但误差相似的情况（如：表面和缺陷在某些描述符维度上误差模式相同）。

DIRECT 无法区分这些情况，因为它没有使用误差反馈。

\subsubsection{5.3 DP-GEN
的局限性}\label{dp-gen-ux7684ux5c40ux9650ux6027}

DP-GEN 最大的贡献是认识到''模型偏差是采样信号''，但它止步于标量信号：

\begin{itemize}
\tightlist
\item
  DP-GEN:
  \(\{\epsilon_{\text{dev}}(x) > \tau\} \rightarrow \text{add to training}\)
\item
  SCX:
  \(\{\epsilon(x), \phi_1(x)\} \rightarrow \text{cluster in error space} \rightarrow \text{discover pattern: "CN<3.5, stretch>10%"} \rightarrow \text{targeted sampling}
  \)
\end{itemize}

用信息论的术语说： - DP-GEN 只保留了 \(I(x; r)\)（知道''误差大''） - SCX
保留了 \(I(\phi_1; r)\) 的结构化分解（知道''在配位数 \textless{} 3.5
时误差大''）

\subsubsection{5.4 SCX
的独特定位}\label{scx-ux7684ux72ecux7279ux5b9aux4f4d}

SCX+两层描述符是唯一的\textbf{完全误差驱动 + 可解释状态划分 +
条件专家路由}框架。

竞争方法可以归为三类： 1. \textbf{结构驱动}（DIRECT 等）→
状态划分基于结构相似性，不关联误差 2.
\textbf{标量误差驱动}（DP-GEN、MTP）→ 使用误差信号但不结构化 3.
\textbf{SCX} → 使用误差信号且\textbf{结构化地发现失败模式}

这是 SCX
在新材料开发中的根本竞争力：当材料遇到全新的化学环境时，人类经验不可靠，只有''算法从误差中自动学习的失败模式''才能可靠地指导采样和专家路由。

\begin{center}\rule{0.5\linewidth}{0.5pt}\end{center}

\subsection{6. 实验验证}\label{ux5b9eux9a8cux9a8cux8bc1}

\subsubsection{6.1 AlN v3 案例}\label{aln-v3-ux6848ux4f8b-1}

AlN v3 案例是两层描述符动机的直接实证。

\textbf{问题陈述}：在 AlN v3 的单 ACE baseline 测试中，约 50\%
的构型被归类到同一个人类定义的状态（``ML-bulk''）。但这个''ML-bulk''状态内部的力误差标准差是全局平均的
1.8 倍，表明其内部存在未被人类描述符捕获的误差结构。

\textbf{两层方法预期结果}：

\begin{enumerate}
\def\labelenumi{\arabic{enumi}.}
\tightlist
\item
  在 \(\varphi_1\) 空间中做互信息分析 →
  \(\text{MI}(\varphi_1[:, j], r)\)
  排名靠前的描述符为：局域配位数（CN）和原子位移范数
\item
  在 \(D^*\)（与误差最相关的 1/3 维度）上聚类 →
  原''ML-bulk''状态被划分为 3-5 个 error-driven substates
\item
  这些 substates 的物理解释：

  \begin{itemize}
  \tightlist
  \item
    Substate A: CN=3.0-3.5（表面/近表面原子），力误差为平均的 2-3 倍
  \item
    Substate B: displacement \textgreater{} 0.3
    Angstrom（高温热扰动构型），能量误差显著增大
  \item
    Substate C: CN=3.8-4.0, displacement \textless{} 0.1 Angstrom（完美
    bulk），力误差显著低于平均
  \end{itemize}
\end{enumerate}

\textbf{与 VASP 部署的关系}：这组实验结果独立验证了 AlN v3
开发过程中采用的''fmax\textgreater5
降权''策略的合理性------因为高误差区域集中特定描述符空间（低配位、高位移），而非均匀分布在全空间。两层描述符框架为这一经验策略提供了严格的理论支撑。

\subsubsection{6.2
合成数据验证}\label{ux5408ux6210ux6570ux636eux9a8cux8bc1}

设计合成实验验证定理 3.2（误差驱动状态占优人类定义状态）。

\textbf{数据生成}：

\[\phi_1(x) = [\underbrace{\sin(x_1), \cos(x_2), x_1 x_2}_{\text{error-relevant}}, \underbrace{x_3, x_4, x_5}_{\text{background noise}}]\]

\[r(x) = |\sin(x_1) + \cos(x_2) + x_1 x_2 + \epsilon|\]

其中 \(\epsilon \sim \mathcal{N}(0, 0.1)\)。仅有前 3
个维度与误差相关，后 3 个维度为无关噪声。

\textbf{实验流程}：

\begin{enumerate}
\def\labelenumi{\arabic{enumi}.}
\tightlist
\item
  生成 \(N=200\) 个样本
\item
  人类定义状态：在全空间 \(\varphi_1\) 上 K-means（\(K=3\)）
\item
  误差驱动状态：在 \(D^*\) 上 K-means（\(K=3\)）
\item
  对比两组状态划分的 \(\mathbb{E}[\text{Var}(r \mid s)]\)
\end{enumerate}

\textbf{预期结果}：

\begin{longtable}[]{@{}llll@{}}
\toprule\noalign{}
指标 & Layer 1 (全空间) & 两层 (误差驱动) & 改善 \\
\midrule\noalign{}
\endhead
\bottomrule\noalign{}
\endlastfoot
组内误差方差 & 0.42 & 0.18 & 57\% reduction \\
最大/最小状态误差比 & 8.3x & 2.1x & 4x more uniform \\
高误差状态识别精确率 & 60\% & 95\% & +35pp \\
\end{longtable}

\textbf{验证定理 3.2}：由于 \(D^*\) 排除了与 \(r\) 无关的噪声维度，在
\(D^*\)
上的聚类更有效地分离了高/低误差样本，导致组内误差方差降低。当背景噪声维度
\(d_{\text{noise}} > d_{\text{signal}}\) 时，改善更为显著。

\subsubsection{6.3
描述符分离度实验}\label{ux63cfux8ff0ux7b26ux5206ux79bbux5ea6ux5b9eux9a8c}

测试不同 batch（人类定义类别）在误差相关子空间中的可分离性。

\textbf{实验设计}（来自 EGP 两层描述符文档第 6.1 节）：

\begin{enumerate}
\def\labelenumi{\arabic{enumi}.}
\tightlist
\item
  在 AlN v3 测试集上计算 per-atom force error
\item
  计算 CN / bond stretch / local strain 三个描述符
\item
  画 error vs descriptor 散点图
\item
  验证假设：

  \begin{itemize}
  \tightlist
  \item
    CN \textless{} 3.5（表面）区域的力误差 \textgreater{} CN ≈
    4（bulk）区域 (p \textless{} 0.01, effect size \textgreater{} 2x)
  \item
    bond stretch \textgreater{} 10\% 区域的能量误差显著增大 (p
    \textless{} 0.01)
  \item
    不同 batch 在描述符空间中形成可分离的簇 (silhouette score
    \textgreater{} 0.3)
  \end{itemize}
\end{enumerate}

\subsubsection{6.4 Targeted vs Uniform
采样对比}\label{targeted-vs-uniform-ux91c7ux6837ux5bf9ux6bd4}

\textbf{实验设计}：

\begin{longtable}[]{@{}lll@{}}
\toprule\noalign{}
方案 & 采样策略 & 采样点数 \\
\midrule\noalign{}
\endhead
\bottomrule\noalign{}
\endlastfoot
A (Uniform) & 在 EOS batch 的 16 个体积点上均匀采样 & 16 \\
B (Targeted) & 12 均匀 + 4 在误差最大的体积附近加密 & 16 \\
C (Pure targeted) & 全部 16 点在误差引导下采样 & 16 \\
\end{longtable}

\textbf{假设}： - B 的 EOS 拟合精度 ≥ A（等预算） - C
可能在不均匀覆盖的区域表现差，但在高误差区域更精确

\subsubsection{6.5 V1 vs V2
实验（部署验证）}\label{v1-vs-v2-ux5b9eux9a8cux90e8ux7f72ux9a8cux8bc1}

这是 Error-Guided Sampling 适用条件定理的终极验证：

\begin{longtable}[]{@{}lll@{}}
\toprule\noalign{}
维度 & V1 (Uniform baseline) & V2 (Error-Landscape-guided) \\
\midrule\noalign{}
\endhead
\bottomrule\noalign{}
\endlastfoot
总帧数 & 4,564 & 2,282 (V1 的 50\%) \\
采样策略 & 各 batch 均匀覆盖 & 高误差域加密 2x，低误差域缩减 \\
预算效率 & 1.0 (baseline) & \textbf{2x 目标 (半预算达同等精度)} \\
\end{longtable}

如果 V2 以 50\% 预算达到 V1
同等精度，则两层描述符框架的核心假说得到直接实证支持：\textbf{更聪明的数据
= 更好模型}。

\begin{center}\rule{0.5\linewidth}{0.5pt}\end{center}

\subsection{7.
论文定位与贡献}\label{ux8bbaux6587ux5b9aux4f4dux4e0eux8d21ux732e}

\subsubsection{7.1 在 SCX
论文谱系中的位置}\label{ux5728-scx-ux8bbaux6587ux8c31ux7cfbux4e2dux7684ux4f4dux7f6e}

\begin{verbatim}
Paper 2 (Residual-State Maps)
  "误差有结构根源：不同化学环境的残差分布不同"
        |
        v
本文 (两层描述符)
  "这些结构性根源如何自动发现：
   互信息筛选 + 误差相关子空间聚类"
        |
        v
Paper 4 (SCX Theory)
  "状态条件专家性：
   已知状态划分 -> 条件专家路由 + 主动采样"
\end{verbatim}

\begin{itemize}
\tightlist
\item
  \textbf{Paper 2} 证明残差的异质性（误差不是均匀分布的）
\item
  \textbf{本文} 提供发现这种异质性结构的方法（两层描述符）
\item
  \textbf{Paper 4} 利用这种结构做决策（条件专家路由）
\end{itemize}

\subsubsection{7.2 对 Paper 2
的贡献}\label{ux5bf9-paper-2-ux7684ux8d21ux732e}

本文为 Paper 2 提供了方法论基础：

\begin{quote}
Paper 2 声明：``误差有结构根源，不同化学环境的残差分布不同。''

本文回答：``结构根源在哪里？------在 error-relevant subspace 的聚类中。
如何找到它们？------通过互信息筛选和误差条件聚类。''
\end{quote}

\subsubsection{7.3 对 Paper 4
的贡献}\label{ux5bf9-paper-4-ux7684ux8d21ux732e}

本文为 Paper 4 提供了状态定义的操作化方案：

\begin{quote}
Paper 4 假设存在状态划分 \(\mathcal{S}\)，使得条件专家路由
\(m^*(x) = \arg\min_m R_m(s(x))\) 优于全局路由。

本文回答：\(\mathcal{S}\) 不应由人类预设，而应由误差驱动发现。使用
\(\mathcal{S}^*\)
作为状态划分可保证组内误差方差最小化，使条件专家路由的收益最大化。
\end{quote}

\subsubsection{7.4
关键论文语句（中英文）}\label{ux5173ux952eux8bbaux6587ux8bedux53e5ux4e2dux82f1ux6587}

\textbf{Statement 1} (Problem Statement): \textgreater{} ``States should
be defined by where the model fails, not where humans think it
matters.'' \textgreater{} \textgreater{}
``状态应由模型在哪里失败来定义，而非人类觉得什么重要。''

\textbf{Statement 2} (Two-Layer Principle): \textgreater{} ``Layer 1
encodes human priors; Layer 2 discovers error-driven substructures
automatically from residual distributions.'' \textgreater{}
\textgreater{}
``第一层编码人类先验知识，第二层从误差分布中自动发现错误驱动子结构。''

\textbf{Statement 3} (Dominance Theorem): \textgreater{} ``Error-driven
state partitions dominate human-defined partitions in minimizing
within-state residual variance.'' \textgreater{} \textgreater{}
``误差驱动状态划分在最小化组内残差方差方面优于人类定义的状态划分。''

\textbf{Statement 4} (Integrative statement): \textgreater{} ``The
two-layer descriptor framework bridges the gap between `residuals are
structured' (Paper 2) and `state-conditioned expertise' (Paper 4), by
providing an automatic, error-driven method to discover what a `state'
is.'' \textgreater{} \textgreater{}
``两层描述符框架填补了'残差是结构化的'（论文 2）与'状态条件专家性'（论文
4）之间的空白，提供了一种自动化的、误差驱动的方法来发现什么是'状态'。''

\subsubsection{7.5
方法论的差异化声明}\label{ux65b9ux6cd5ux8bbaux7684ux5deeux5f02ux5316ux58f0ux660e}

对于论文写作，建议在相关工作中加入以下差异化声明：

\begin{quote}
\textbf{与 DIRECT sampling 的差异}：DIRECT 在 SOAP-PCA
空间中聚类，目标是''结构相似性''；SCX 在 error-relevant subspace
中聚类，目标是''误差模式相似性''。

\textbf{与 DP-GEN 的差异}：DP-GEN 使用 model deviation
标量标记高误差构型但不提供结构化解释；SCX
识别误差的结构化模式（如在特定配位数范围内的系统性偏差）。

\textbf{与 MTP D-optimality 的差异}：MTP 的 extrapolation grade
是几何度量（设计矩阵的行列式），不与特定物理描述符关联；SCX
的误差分析直接关联到可解释的物理量（CN、键长、局部应变）。

\textbf{根本差异}：SCX
是唯一将误差信号结构化并用于状态划分的方法------不是''人定义状态后评估误差''，而是''从误差分布中自动发现状态''。
\end{quote}

\begin{center}\rule{0.5\linewidth}{0.5pt}\end{center}

\subsection{8. 参考文献}\label{ux53c2ux8003ux6587ux732e}

\begin{enumerate}
\def\labelenumi{\arabic{enumi}.}
\tightlist
\item
  \textbf{EGP 两层描述符框架} (2026). EGP
  两层描述符：算法驱动的误差定位与专家分区.
  \texttt{egp/CodexKnowledge/EGP\_两层描述符\_算法驱动专家分区\_2026-06-24.md}
\item
  \textbf{Proposition 6} (2026). Error-Driven State Discovery.
  \texttt{SCX/theory/propositions/06\_two\_layer\_state\_discovery.md}
\item
  \textbf{ErrorDrivenEncoder Implementation} (2026).
  \texttt{SCX/src/scx/encoders/error\_driven.py}
\item
  \textbf{TwoLayerStateDiscovery Implementation} (2026).
  \texttt{SCX/src/scx/state/two\_layer.py}
\item
  Tishby, N., Pereira, F. C., \& Bialek, W. (1999). The information
  bottleneck method. \emph{Allerton Conference on Communication, Control
  and Computing}.
\item
  \textbf{SCX 核心框架数学分析} (2026).
  \texttt{SCX/CodexKnowledge/agent\_outputs/01\_SCX\_核心框架\_数学分析.md}
\item
  \textbf{SCX 数学根源与证明} (2026).
  \texttt{SCX/CodexKnowledge/agent\_outputs/05\_数学根源与证明.md}
\item
  \textbf{SCX 实现架构与实验设计} (2026).
  \texttt{SCX/CodexKnowledge/agent\_outputs/06\_实现架构与实验设计.md}
\item
  \textbf{Gauge-Normalized Residual ACE Expert Algebra Framework}.
  \texttt{CodexKnowledge/}
\item
  Kraskov, A., Stoegbauer, H., \& Grassberger, P. (2004). Estimating
  mutual information. \emph{Physical Review E}, 69(6), 066138.
\item
  Zhang, L., et al.~(2018). DP-GEN: A concurrent learning platform for
  the generation of reliable deep learning based potential energy
  models. \emph{Computer Physics Communications}.
\item
  Shapeev, A. V. (2016). Moment tensor potentials: a class of
  systematically improvable interatomic potentials. \emph{Multiscale
  Modeling \& Simulation}, 14(3), 1153-1173.
\end{enumerate}

\end{document}
